\section{Conclusion}
% =========================================================

This work presents an implemented preliminary-design framework for a single bell
nozzle. A continuous piecewise contour composed of a straight convergent, circular
subsonic and supersonic throat arcs, and a quadratic divergent section supports
two-dimensional plotting, axisymmetric three-dimensional visualization, CAD export
and quasi-one-dimensional flow reconstruction. The throat radius is retained from
engine pre-sizing, while the exit radius follows from the expansion ratio.

The genetic algorithm varies three real-valued design variables:
\(\varepsilon\), \(K_b\) and \(\theta_{in}\). The convergent angle and reference
cone half-angle remain user-selected fixed inputs, and the exit wall angle is derived
from the resulting quadratic contour. The optimized objective is the effective
ambient thrust coefficient, combining RocketCEA momentum performance, exit divergence,
compressible boundary-layer blockage, signed pressure thrust and integrated wall
friction. No independent curvature or empirical turning efficiency is used.

Two viscous evaluation levels provide a practical cost--fidelity trade-off. Quick
screening uses a weak momentum-integral closure, whereas BLIMP-lite resolves a
wall-normal velocity profile with a Cebeci--Smith eddy-viscosity model and is used for
final reporting. Both assume an attached, fully turbulent and adiabatic wall; hence
\(T_w=T_{aw}\) and the present model does not predict non-zero wall heat flux.

For the representative saved run, the selected geometry reduced total length by
\(17.32\%\) and increased \(C_{F,eff}\) by \(0.41\%\) relative to the baseline.
This result demonstrates the coupled workflow rather than final nozzle qualification.
The framework is therefore suitable for transparent preliminary comparison and
design-space exploration, provided that final candidates are validated using
independent higher-fidelity analysis and experimental evidence.
